% authority: science_draft
% object-id: TEX-V22-P4-T02-B2-COMMON-HYPERBOLICITY-ENVELOPE-CANDIDATE-V1
% task-id: RT-20260810-005
% agentjob-id: AJ-RT-20260810-005-001
\documentclass[11pt]{article}

\usepackage[margin=0.78in]{geometry}
\usepackage{amsmath,amssymb,amsthm}
\usepackage{booktabs,longtable,array,enumitem}
\usepackage[hidelinks]{hyperref}

\newtheorem{definition}{Definition}
\newtheorem{theorem}{Theorem}
\newtheorem{proposition}{Proposition}
\newtheorem{lemma}{Lemma}
\newenvironment{nonconclusion}{\par\smallskip\noindent\textbf{Non-conclusion.}\ }{\par\smallskip}

\newcommand{\RR}{\mathbb R}
\newcommand{\Seq}{\mathcal S_{\mathrm{eq}}}
\newcommand{\Pprod}{P_{\Pi}}
\newcommand{\GammaInt}{\Gamma_{\cap}}
\newcommand{\Resp}{\operatorname{Resp}_{\mathrm{prin}}}
\newcommand{\Qray}{Q_{\mathrm{ray}}}
\newcommand{\Qsign}{Q_{\mathrm{sign}}}

\title{V22 P4--T02 B2 Common-Hyperbolicity Envelope Candidate v1\\
\large A proposal-only product-cone and two-layer response-quotient construction}
\author{Candidate Constructor Packet RT-20260810-005}
\date{2026-08-10}

\begin{document}
\sloppy
\maketitle

\begin{abstract}
This draft/control packet executes the post-line-lock construction selected by
RT-20260810-002.  On the exactly equipped finite sector family
\(\Seq=\{R,S,D\}\), it supplies three independently split, proposal-only
degree-one principal forms, proves hyperbolicity of their cubic product,
identifies the containing component as their common oriented cone, and gives a
source-presentation-covariant strict-margin certificate under independent
sector perturbations.  It also constructs two explicit response quotients:
a positive-projective ray for fixed sector representatives that retains
relative response ratios, and a coarser normalization-invariant
positive-diagonal sign orbit.  The all-positive cell of either quotient has
the common cone as its exact preimage.  These quotients are source-extension
candidates, not implementations of the unchanged P7 protocols and not
physical causal quotients.  Empty-intersection, orientation, zero-factor,
multiplicity, normalization, inverse, cocycle, scope, and target-import
controls fail closed.  The decisive result is one \emph{constructed
candidate}, followed by Smuggling Auditor review.  No all-sector universality,
strong hyperbolicity, physical cone, conformal structure, effective metric,
adequacy verdict, B2 activation, P4--T03 authority, ontology adoption, or
Distance-to-GR reduction follows.
\end{abstract}

\section{Fixed authority and source basis}

The packet is bound to
\[
 \texttt{CAND-V22-B2-COMMON-HYPERBOLICITY-ENVELOPE-V1}
\]
and work item
\[
 \texttt{P4-T02-B2-COMMON-HYPERBOLICITY-ENVELOPE-CONSTRUCTION}.
\]
It preserves the scoped obstruction
\(\texttt{OBST-V22-P4T02-B2-NATURAL-LINE-LOCK-001}\) and the local
\(\texttt{NDCL-V22-P4T02-B2-SHARED-TAU-SELECTOR}\) freeze.  Consequently,
no common sector line is inserted and no diagonal-only variation class is
used.

The exact positive domain contains only the equipped
\(\mathsf{Rod}_{\rm src}\), \(\mathsf{Signal}_{\rm src}\), and
\(\mathsf{Detector}_{\rm src}\) controls.  Their existing P7 protocols remain
unchanged.  The smooth patch, principal forms, orientation, and new
principal-response request/readout rule below are proposal-only source-extension
data.  Current ontology does not derive them.  No target atlas, metric,
tetrad, coframe, desired GR cone, empirical calibration, or benchmark answer
is an input.

\section{Intrinsic finite construction}

Let \(U\) be a declared smooth four-dimensional source patch, let
\(x\in U\), and write \(V=T_xU\).  Supply, as proposal-only data, a covector
\(h_0\in V^*\), vectors \(u,z\in V\), and
\(\epsilon=1/3\) satisfying
\[
 h_0(u)=1,\qquad h_0(z)=0,\qquad z\ne0.
\]
Define the independently split sector vectors
\[
 v_R=u+\epsilon z,\qquad v_S=u,\qquad v_D=u-\epsilon z
\]
and the degree-one principal forms
\[
 p_s(k)=k(v_s),\qquad s\in\Seq,\quad k\in V^*.
\]
These forms have density weight zero in the candidate.  They are nonzero and
are not associates when \(\epsilon\ne0\).  Their orientation receipt is
\[
 p_R(h_0)=p_S(h_0)=p_D(h_0)=1>0.
\]

\begin{definition}[Finite common-hyperbolicity envelope]
Define
\begin{align}
 \Pprod(k)
   &=p_R(k)p_S(k)p_D(k) \\
   &=k(u)\bigl(k(u)^2-\epsilon^2 k(z)^2\bigr), \\
 \GammaInt
   &=\{k\in V^*:p_R(k)>0,\ p_S(k)>0,\ p_D(k)>0\} \\
   &=\{k\in V^*:k(u)>|\epsilon k(z)|\}.
\end{align}
The pair \((\Pprod,\GammaInt)\) is the finite mathematical envelope.
\end{definition}

\begin{theorem}[Constructed product and common component]
\label{thm:product}
The polynomial \(\Pprod\) is hyperbolic with respect to \(h_0\).  Its
hyperbolicity-cone component containing \(h_0\) is exactly \(\GammaInt\).
The cone is nonempty, open, and convex, while the three characteristic
hyperplanes remain distinct.
\end{theorem}

\begin{proof}
For every \(q\in V^*\), the roots of
\(t\mapsto\Pprod(q+t h_0)\) are
\[
 -q(v_R),\qquad -q(v_S),\qquad -q(v_D),
\]
because \(h_0(v_s)=1\).  All roots are real, proving hyperbolicity.  The
positive component containing \(h_0\) is the simultaneous positivity region
of the three factors, hence \(\GammaInt\).  The displayed inequality makes
openness and convexity explicit, and \(h_0\in\GammaInt\).  Since
\(z\ne0\) and \(\epsilon\ne0\), the three linear forms do not have one common
kernel, so their characteristic hyperplanes are distinct.
\end{proof}

The theorem is the exact local mathematical payload.  It neither says that
\(\Pprod\) is the physical principal polynomial of each sector nor converts
the union of three characteristic hyperplanes into one shared boundary.

\section{Source-presentation covariance and gluing obligations}

Let \(A:V\to V'\) be an invertible derivative of an admitted source
presentation change.  Transport the proposal-only data by
\[
 v'_s=A v_s,\qquad k=A^*k'\quad(k'\in V'^*).
\]
Then
\[
 p'_s(k')=k'(A v_s)=p_s(A^*k'),
\]
so products and common components obey
\[
 \Pprod'(k')=\Pprod(A^*k'),
 \qquad
 \GammaInt'=(A^*)^{-1}(\GammaInt).
\]
This is source covariance conditional on the supplied data; it is not a
derivation of those data from a bare smooth patch.

On overlaps, the derivatives must be invertible and satisfy the ordinary
triple-overlap cocycle.  Sector relabelings act by simultaneous permutation of
the three forms.  A common positive rescaling of the three fixed sector
representatives is absorbed by the projective ray quotient below.  Independent
positive generator rescalings generally change its response ratios, but are
absorbed by the coarser positive-diagonal sign quotient.  A vanishing or
sign-changing transition factor fails the oriented gluing contract.  A failed
inverse or triple-overlap cocycle blocks a global envelope rather than being
repaired by target geometry.

\section{Operational two-layer response-quotient candidate}

The exact P7 protocols currently contain neither a covector request nor a
principal-response score.  The following law is therefore a new proposal-only
protocol extension, not a statement about current devices.

On the domain
\[
 D=\{k\in V^*:p_s(k)\ne0\ \text{for every }s\in\Seq\},
\]
define the formal sector response vector
\[
 \Resp(k)=\bigl(p_R(k),p_S(k),p_D(k)\bigr)\in(\RR\setminus\{0\})^3.
\]
For the fixed proposal-only representatives above, first quotient by common
positive scale:
\[
 \Qray(k)=[\Resp(k)]_{\RR_{>0}}\in
 \bigl((\RR\setminus\{0\})^3\bigr)/\RR_{>0}.
\]
This response ray retains the two independent relative response ratios.  It is
invariant under a common positive rescaling of all three generators, but it
is not invariant under independent sector rescalings.  Thus it is a legitimate
local proposal-only classifier for the explicitly fixed representatives; a
canonical global ray would additionally require a coherent relative-sector
normalization law.

For the normalization-insensitive layer, let
\(G_+=(\RR_{>0})^3\) act componentwise and define
\[
 \Qsign(k)=[\Resp(k)]_{G_+}.
\]
Equivalently, \(k\sim_{\rm sign}k'\) when there are
\(a_R,a_S,a_D>0\) such that
\(p_s(k')=a_s p_s(k)\) for every sector.  This coarser quotient retains only
the ordered sign pattern and discards all independent positive generator
normalizations.

\begin{theorem}[Two all-positive preimages]
\label{thm:quotient}
Let \(\mathcal C_{+++}\) be the positive projective cell and let
\(\mathcal O_{+++}=G_+\cdot(1,1,1)\).  Then
\[
 \GammaInt=\Qray^{-1}(\mathcal C_{+++})
           =\Qsign^{-1}(\mathcal O_{+++}).
\]
The ray identity is invariant under common positive scale; the sign-orbit
identity is invariant under independent positive local generator rescaling.
Both are covariant under admitted source-presentation transport and
simultaneous sector relabeling.
\end{theorem}

\begin{proof}
Membership in either \(\mathcal C_{+++}\) or \(\mathcal O_{+++}\) is
equivalent to positivity of all three components of \(\Resp(k)\), which is
precisely the definition of \(\GammaInt\).  Common and positive-diagonal
rescaling preserve component signs.  Source-presentation covariance preserves
each evaluation, while relabeling permutes the components and the respective
all-positive subset.
\end{proof}

The proposed protocol law adds a source-covector preparation label and returns
the formal score \(p_s(k)\), from which one of the tokens
\(\mathsf{principal\_positive}\),
\(\mathsf{principal\_zero}\), or
\(\mathsf{principal\_negative}\) is derived according to its sign.  The ray layer
discards one common positive scale while retaining ratios; the sign layer
discards independent positive normalizations.  The rule is noncircular with
respect to the target because it uses only source data and evaluation signs.
It is nonetheless an unadopted new primitive: no existing P7 response theorem
makes these formal signs or ratios empirical detector responses or physical
signal outcomes.  In particular, the ray layer does not itself supply the
missing relative-normalization law, and the sign layer does not retain the
information that such a law would make meaningful.

\section{Strict-margin independent-variation certificate}

\begin{proposition}[Intrinsic finite strict margin]
There are product-open neighborhoods \(N_h\) of \(h_0\) and \(N_s\) of each
\(v_s\) such that
\[
 h'(v'_s)>0
\]
for every independently selected
\((h',v'_R,v'_S,v'_D)\in N_h\times N_R\times N_S\times N_D\).
Consequently the perturbed common intersection is nonempty and contains
\(h'\).
\end{proposition}

\begin{proof}
The evaluation pairing \(V^*\times V\to\RR\) is continuous.  Every base
evaluation equals one, and the family is finite.  Intersecting the finitely
many inverse images of \((0,\infty)\) gives the required product-open
neighborhood.  The construction does not constrain the sector variations to
be diagonal.
\end{proof}

For reproducibility only, take
\(u=e_0\), \(z=e_1\), and \(h_0=e^0\).  If every coordinate component of
\(h'-h_0\) and \(v'_s-v_s\) has magnitude at most
\(\delta=1/10\), then
\[
 |h'(v'_s)-1|
 \le \frac43\delta+\delta+4\delta^2
 =\frac{41}{150},
\]
and hence \(h'(v'_s)\ge109/150>0\).  This numeric sup bound is a chart-level
receipt for the intrinsic openness proof.  It is not a source norm, a target
metric, or a global conditioning theorem.

\section{Adversarial controls and exact limitations}

The common-direction premise is nontrivial.  At one point let
\[
 w_1=e_0,\qquad w_2=e_1,\qquad w_3=-e_0-e_1.
\]
Each linear form \(k\mapsto k(w_i)\) is individually hyperbolic, but no
covector can make all three positive: the first two inequalities imply the
negative of the third.  Equivalently, \(w_1+w_2+w_3=0\) is a positive
dependence certificate.  This negative control fails closed with empty
intersection.

Erasing the orientation token leaves the all-positive and all-negative
components equally available.  Setting one sector vector to zero collapses
the product.  Singular overlap maps, vanishing transition factors, or failed
cocycles block globalization.  Requesting an unequipped P7 sector is outside
the declared domain rather than evidence for universality.

The cubic product is reducible and not strictly hyperbolic on all covectors:
sector roots coincide when, for example, \(q(z)=0\).  The construction does
not establish strong or symmetric hyperbolicity, multiplicity stability,
energy estimates, finite condition numbers, one signal speed, or one physical
characteristic boundary.  These are surviving P4--T02 burdens.

\section{Source-extension classification and next role}

Finite product, finite cone intersection, common-scale projectivization,
positive-diagonal orbit formation, and pullback covariance are conservative
mathematical operations once the inputs are supplied.  The sector principal
forms, orientation, relative-normalization status, and
protocol-to-principal-response law are a
\emph{new ontology primitive candidate}, with status
\emph{proposal-only} and material class \emph{source-extension data}.
Current ontology does not derive the
\texttt{source\_operational\_common\_hyperbolicity\_envelope\_quotient\_law}.
No adoption is requested.  Current adoption is blocked while same-milestone
research continuation remains open:
\texttt{blocked\_adoption\_open\_continuation}.

The next role is \texttt{smuggling-auditor@0.2.0}.  It must test whether the
new covector request, response scores, sign tokens, orientation, positive
transition factors, sector restriction, and operational vocabulary hide target
or goal-property imports.  The audit is not executed in this task.

\section{Distance-to-GR status}

\begin{longtable}{>{\raggedright\arraybackslash}p{0.24\textwidth}>{\raggedright\arraybackslash}p{0.23\textwidth}>{\raggedright\arraybackslash}p{0.43\textwidth}}
\toprule
Burden & Tracked status preserved & Task-local consequence\\
\midrule
Source ontology primitives & accepted in scoped source-extension status & Adds unadopted envelope inputs; no canonical source law changes.\\
Source equivalence EqSrc & draft object exists & No general equivalence theorem is added.\\
RetainH & blocked by missing primitive & Unchanged.\\
GenH & blocked by missing primitive & Unchanged.\\
ObsLoc\(_{\rm lc}\) & constructive witness exists & Exact equipped-protocol scope is read only; no new localization semantics.\\
Resp\(_{\rm lc}\) & accepted only as scoped source-extension data & The principal-response law is a new unadopted proposal.\\
\(M_{\rm src}\) & accepted only in its scoped protected status & The smooth patch remains conditional input, not a target manifold.\\
\(g_{\rm eff}\) & accepted only as the historical scoped source-extension record; bare burden unresolved & The envelope is not a metric or physical cone.\\
Matter coupling & accepted only by the historical scoped source-side postulate & No target coupling or universal propagation is derived.\\
Einstein equations & draft object exists with positive derivation blocked & No action, variation, correction control, or field equation is supplied.\\
Finite-variation robustness & Refuter stress passed in the tracked ledger & Adds one local independent-sector open-neighborhood witness without changing the ledger.\\
Benchmark promotion & blocked by missing primitive & No benchmark comparison or promotion.\\
Gate Chair status & human-gated & No protected review is requested.\\
Current route freeze or hard-fail status & prior shared-line route locally frozen & The materially distinct envelope candidate is constructed and awaits audit; no broader freeze is inferred.\\
\bottomrule
\end{longtable}

The ledger is not modified.  The task contributes new mathematical payload but
does not reduce the governed Distance-to-GR burden.

\section{No-fog result}

Exactly one decisive result is recorded:
\begin{quote}
\texttt{constructed\_candidate} --- the exact finite local product-cone
envelope, fixed-representative response ray, normalization-invariant sign
orbit, covariance map, strict-margin certificate, and adversarial controls
exist at the declared proposal-only source-extension level.  The next
required role is Smuggling Auditor.
\end{quote}

The constructed object is useful because it proves that a common time-covector
domain can be defined robustly without equality of sector characteristic
lines.  Its limitation is equally decisive: the new inputs and response
law are not derived or adopted, the exact P7 protocols do not implement them,
and distinct sector boundaries plus unresolved multiplicities remain.

\begin{nonconclusion}
The candidate is not a physical cone, universal matter propagation law,
Lorentzian conformal structure, time orientation for physics, metric, scale,
matter coupling, or Einstein equation.  The three equipped controls do not
exhaust P7.  D7 adequacy is unevaluated, B2 is inactive, P4--T03 is locked,
the Distance-to-GR ledger is unchanged, and no ontology adoption, Gate
verdict, benchmark promotion, proof authority, publication, push, external
action, global no-go, future-extension impossibility, or completed derivation
follows.
\end{nonconclusion}

\end{document}
